% GENERATED FILE. Edit canonical lessons, then run npm run generate:books.
% canonical-source-hash: fnv1a64:83de6bde07bf6f1c
% canonical-lessons: GU-C484-ilaj, GU-C484-upcar, GU-C484-nidan, GU-C484-matadvu, GU-C484-injeksan

\chapter{Finding a Cure}
\label{ch:gu-finding-a-cure}
\hlchaptermodality{484}

\begin{chapteropening}
\textbf{By the end of this chapter:} \emph{I can say how a doctor treats an illness.}

Complete the last lesson of chapter 484: I can say how a doctor treats an illness.
\end{chapteropening}

\section[\texorpdfstring{\gu{ઇલાજ} (ilāj)}{ઇલાજ (ilāj)}]{\gu{ઇલાજ} (ilāj) — a cure}

\label{lesson:GU-C484-ilaj}



\begin{quote}
Before the new one: say the Gujarati for to writhe, then the Gujarati for a burning feeling.
\end{quote}

\subsection*{You'll want to know: \gu{ઇલાજ}}
\textbf{\gu{ઇલાજ}} — \emph{ilāj} — \textquotedblleft{}a cure\textquotedblright{}.

\begin{quote}
\gu{આ} \gu{રોગનો} \gu{ઇલાજ} \gu{છે}, \gu{ચિંતા} \gu{ન} \gu{કરો}.
\end{quote}

\begin{quote}
\emph{ā rogno ilāj che, cintā na karo.}
\end{quote}

\begin{quote}
There is a cure for this illness; don't worry.
\end{quote}

A masculine noun.

\begin{cousinweb}
\textbf{\gu{ઇલાજ}} (\emph{ilāj}) is usually traced to Arabic \emph{ʿilāj}, \textquotedblleft{}treatment, cure\textquotedblright{}, which reached Gujarati through Persian.
\end{cousinweb}

\begin{grammarlens}[title={What you've built}]
One more word for describing what you lived and felt.
\end{grammarlens}

\subsection*{Your turn}
\begin{itemize}
\raggedright
  \item \emph{Say it:} \emph{ilāj}
  \item \emph{Say it:} \emph{ilāj}, once more
  \item \emph{Say it:} \emph{ā rogno ilāj che, cintā na karo.} — There is a cure for this illness; don't worry.
  \item \emph{Recall:} say the Gujarati for to writhe, then the Gujarati for a burning feeling, then say \emph{ilāj} again
\end{itemize}

\begin{tcolorbox}[breakable,colback=teal!4,colframe=teal!35!black,title={Before you move on}]
What does \gu{ઇલાજ} mean? (\textquotedblleft{}A cure\textquotedblright{}.) Say it once more.
\end{tcolorbox}
\section[\texorpdfstring{\gu{ઉપચાર} (upcār)}{ઉપચાર (upcār)}]{\gu{ઉપચાર} (upcār) — a remedy}

\label{lesson:GU-C484-upcar}



\begin{quote}
Before the new one: say the Gujarati for a burning feeling, then the Gujarati for a cure.
\end{quote}

\subsection*{You'll want to know: \gu{ઉપચાર}}
\textbf{\gu{ઉપચાર}} — \emph{upcār} — \textquotedblleft{}a remedy\textquotedblright{}.

\begin{quote}
\gu{દાદી} \gu{શરદી} \gu{માટે} \gu{આદુનો} \gu{ઉપચાર} \gu{કરે} \gu{છે}.
\end{quote}

\begin{quote}
\emph{dādī śardī māṭe āduno upcār kare che.}
\end{quote}

\begin{quote}
Grandmother treats a cold with ginger.
\end{quote}

A masculine noun.

\begin{cousinweb}
\textbf{\gu{ઉપચાર}} (\emph{upcār}) is a tatsama, Sanskrit \emph{upacāra}, \textquotedblleft{}approach, service, attendance\textquotedblright{}, as Monier-Williams records it; the medical sense grew from attending on the sick.
\end{cousinweb}

\begin{grammarlens}[title={What you've built}]
One more word for describing what you lived and felt.
\end{grammarlens}

\subsection*{Your turn}
\begin{itemize}
\raggedright
  \item \emph{Say it:} \emph{upcār}
  \item \emph{Say it:} \emph{upcār}, once more
  \item \emph{Say it:} \emph{dādī śardī māṭe āduno upcār kare che.} — Grandmother treats a cold with ginger.
  \item \emph{Recall:} say the Gujarati for a burning feeling, then the Gujarati for a cure, then say \emph{upcār} again
\end{itemize}

\begin{tcolorbox}[breakable,colback=teal!4,colframe=teal!35!black,title={Before you move on}]
What does \gu{ઉપચાર} mean? (\textquotedblleft{}A remedy\textquotedblright{}.) Say it once more.
\end{tcolorbox}
\section[\texorpdfstring{\gu{નિદાન} (nidān)}{નિદાન (nidān)}]{\gu{નિદાન} (nidān) — a diagnosis}

\label{lesson:GU-C484-nidan}



\begin{quote}
Before the new one: say the Gujarati for a cure, then the Gujarati for a remedy.
\end{quote}

\subsection*{You'll want to know: \gu{નિદાન}}
\textbf{\gu{નિદાન}} — \emph{nidān} — \textquotedblleft{}a diagnosis\textquotedblright{}.

\begin{quote}
\gu{દાક્તરે} \gu{રોગનું} \gu{સાચું} \gu{નિદાન} \gu{કર્યું}.
\end{quote}

\begin{quote}
\emph{dāktare rognũ sācũ nidān karyũ.}
\end{quote}

\begin{quote}
The doctor made the right diagnosis.
\end{quote}

A neuter noun.

\begin{cousinweb}
\textbf{\gu{નિદાન}} (\emph{nidān}) is a tatsama, Sanskrit \emph{nidāna}, which Monier-Williams gives as \textquotedblleft{}a band\textquotedblright{} and as a first cause; the old medicine used it for finding the cause of an illness.
\end{cousinweb}

\begin{grammarlens}[title={What you've built}]
One more word for describing what you lived and felt.
\end{grammarlens}

\subsection*{Your turn}
\begin{itemize}
\raggedright
  \item \emph{Say it:} \emph{nidān}
  \item \emph{Say it:} \emph{nidān}, once more
  \item \emph{Say it:} \emph{dāktare rognũ sācũ nidān karyũ.} — The doctor made the right diagnosis.
  \item \emph{Recall:} say the Gujarati for a cure, then the Gujarati for a remedy, then say \emph{nidān} again
\end{itemize}

\begin{tcolorbox}[breakable,colback=teal!4,colframe=teal!35!black,title={Before you move on}]
What does \gu{નિદાન} mean? (\textquotedblleft{}A diagnosis\textquotedblright{}.) Say it once more.
\end{tcolorbox}
\section[\texorpdfstring{\gu{મટાડવું} (maṭāḍvũ)}{મટાડવું (maṭāḍvũ)}]{\gu{મટાડવું} (maṭāḍvũ) — to cure}

\label{lesson:GU-C484-matadvu}



\begin{quote}
Before the new one: say the Gujarati for a remedy, then the Gujarati for a diagnosis.
\end{quote}

\subsection*{You'll want to know: \gu{મટાડવું}}
\textbf{\gu{મટાડવું}} — \emph{maṭāḍvũ} — \textquotedblleft{}to cure\textquotedblright{}.

\begin{quote}
\gu{આ} \gu{દવા} \gu{તાવ} \gu{મટાડે} \gu{છે}.
\end{quote}

\begin{quote}
\emph{ā davā tāv maṭāḍe che.}
\end{quote}

\begin{quote}
This medicine cures fever.
\end{quote}

Past, with \gu{તેણે} (he or she …): \gu{તેણે} \gu{રોગ} \gu{મટાડ્યો}, \emph{teṇe rog maṭāḍyo}.

\begin{cousinweb}
\textbf{\gu{મટાડવું}} (\emph{maṭāḍvũ}) is the causative of \textbf{\gu{મટવું}} (\emph{maṭvũ}), \textquotedblleft{}to be cured\textquotedblright{}, which CDIAL lists, \textquotedblleft{}to be rubbed out\textquotedblright{}, under Sanskrit \emph{mṛṣṭa}, \textquotedblleft{}rubbed\textquotedblright{} (CDIAL 10299).
\end{cousinweb}

\begin{grammarlens}[title={What you've built}]
One more word for describing what you lived and felt.
\end{grammarlens}

\subsection*{Your turn}
\begin{itemize}
\raggedright
  \item \emph{Say it:} \emph{maṭāḍvũ}
  \item \emph{Say it:} \emph{maṭāḍvũ}, once more
  \item \emph{Say it:} \emph{ā davā tāv maṭāḍe che.} — This medicine cures fever.
  \item \emph{Recall:} say the Gujarati for a remedy, then the Gujarati for a diagnosis, then say \emph{maṭāḍvũ} again
\end{itemize}

\begin{tcolorbox}[breakable,colback=teal!4,colframe=teal!35!black,title={Before you move on}]
What does \gu{મટાડવું} mean? (\textquotedblleft{}To cure\textquotedblright{}.) Say it once more.
\end{tcolorbox}
\section[\texorpdfstring{\gu{ઇન્જેક્શન} (injekśan)}{ઇન્જેક્શન (injekśan)}]{\gu{ઇન્જેક્શન} (injekśan) — an injection}

\label{lesson:GU-C484-injeksan}



\begin{quote}
Before the new one: say the Gujarati for a diagnosis, then the Gujarati for to cure.
\end{quote}

\subsection*{You'll want to know: \gu{ઇન્જેક્શન}}
\textbf{\gu{ઇન્જેક્શન}} — \emph{injekśan} — \textquotedblleft{}an injection\textquotedblright{}.

\begin{quote}
\gu{નર્સે} \gu{મને} \gu{ઇન્જેક્શન} \gu{આપ્યું}.
\end{quote}

\begin{quote}
\emph{narse mane injekśan āpyũ.}
\end{quote}

\begin{quote}
The nurse gave me an injection.
\end{quote}

A neuter noun.

\begin{cousinweb}
\textbf{\gu{ઇન્જેક્શન}} (\emph{injekśan}) is borrowed from English \emph{injection}.
\end{cousinweb}

\begin{grammarlens}[title={What you've built}]
One more word for describing what you lived and felt.
\end{grammarlens}

\subsection*{Your turn}
\begin{itemize}
\raggedright
  \item \emph{Say it:} \emph{injekśan}
  \item \emph{Say it:} \emph{injekśan}, once more
  \item \emph{Say it:} \emph{narse mane injekśan āpyũ.} — The nurse gave me an injection.
  \item \emph{Recall:} say the Gujarati for a diagnosis, then the Gujarati for to cure, then say \emph{injekśan} again
\end{itemize}

\begin{tcolorbox}[breakable,colback=teal!4,colframe=teal!35!black,title={Before you move on}]
What does \gu{ઇન્જેક્શન} mean? (\textquotedblleft{}An injection\textquotedblright{}.) Say it once more.
\end{tcolorbox}
